\documentclass[10pt]{ctexart}
\usepackage[a4paper,margin=1in]{geometry}
\usepackage{amsmath,amssymb,mathtools,bm}
\usepackage{booktabs}
\usepackage{enumitem}
\usepackage{microtype}
\usepackage{hyperref}
\setlist[itemize]{leftmargin=1.5em}
\newcommand{\clip}{\operatorname{clip}}
\newcommand{\Norm}{\operatorname{Norm}}
\title{001：Step-wise PPO 初步实验规划}
\author{}
\date{}
\begin{document}
\maketitle
\vspace{-1.2em}

\section*{范围}
本文记录 Step-wise PPO 的第一版实验规划。目标是在 \texttt{verl-agent} 的多步 agent 任务中，把训练信号从“每个 step 使用同一个 episode-level reward/advantage”推进到“在环境 step 维度上计算 return、TD error 和 GAE”。这里的 step 指一次 agent 决策：agent 接收当前 context/observation，生成一个完整 action，例如 tool call 或文本动作。

\section*{当前 PPO 实现理解}
标准 PPO 入口为 \path{verl/trainer/main_ppo.py}。训练流程包括：多轮环境交互采样、重新计算 old log probability、可选计算 reference log probability 和 critic value、计算 reward 与 advantage、最后更新 critic 和 actor。

在现有 GAE 路径中，critic value 和 advantage 是 token-shaped：
\[
\delta_\ell = r_\ell + \gamma V_{\ell+1} - V_\ell,
\qquad
A_\ell = \delta_\ell + \gamma\lambda A_{\ell+1}.
\]
这对普通文本 RLHF 是自然的，但对 agentic action 不够严格，因为一个环境动作可能由多个 token 组成，而环境只在 action 粒度返回奖励。

\section*{当前 GRPO 实现理解}
GRPO 不维护 critic。\path{env.rollout.n} 会为同一个初始任务采样多条 trajectory，这些 trajectory 共享同一个 \path{uid}，但各自有不同的 \path{traj_uid}。GRPO 使用 episode return 做 group-relative normalization：
\[
A_i^{\mathrm{GRPO}}
= \frac{R_i - \mu_g}{\sigma_g + \epsilon},
\qquad
\mu_g = \frac{1}{|g|}\sum_{j\in g} R_j.
\]
这个 scalar advantage 会广播到对应 response 的有效 token 上。GRPO 是干净的 non-step baseline，因为它不使用 step reward、step order 或 anchor observation。

\section*{当前 GiGPO 实现理解}
GiGPO 仍然是 critic-free，但它显式利用 step 信息。它先沿 trajectory 反向计算 discounted step return：
\[
G_t = r_t + \gamma G_{t+1},
\]
然后在同一 task group 内按 anchor observation 构造 step group，并在 step group 内做相对归一化。最终 advantage 可以写成：
\[
A_t^{\mathrm{GiGPO}}
= A_i^{\mathrm{epi}} + w_{\mathrm{step}} A_t^{\mathrm{step-group}}.
\]
因此，GiGPO 是评估 StepPPO 时最重要的 step-aware baseline。

\section*{Step-wise PPO 的严格定义}
一条 trajectory 表示为
\[
\tau=(s_0,a_0,r_0,s_1,a_1,r_1,\ldots,s_{T-1},a_{T-1},r_{T-1}).
\]
虽然 action $a_t=(y_{t,1},\ldots,y_{t,L_t})$ 由多个 token 构成，但 value、return、TD error 和 GAE 都在环境 step 维度上定义。第一版采用严格 state-value：
\[
V_t = V(s_t),
\]
默认取 action 生成前最后一个 context token 的 hidden state 作为 value head 输入。取第一个 response token 或最后一个 response token 会更接近 action-conditioned value，应作为 ablation 记录。

Step-wise TD 和 GAE 为
\[
\delta_t = r_t + \gamma(1-d_t)V^{\mathrm{old}}(s_{t+1}) - V^{\mathrm{old}}(s_t),
\]
\[
A_t^{\mathrm{step}} = \delta_t + \gamma\lambda(1-d_t)A_{t+1}^{\mathrm{step}},
\qquad
R_t^{\mathrm{step}} = A_t^{\mathrm{step}} + V^{\mathrm{old}}(s_t).
\]

\section*{Advantage 混合路线}
第一版保留 GRPO/GiGPO 的 episode-level relative signal，同时加入 learned step-wise critic correction：
\[
A_t^{\mathrm{final}}
= \frac{\widehat{A}_t^{\mathrm{epi}} + w_{\mathrm{step}}\widehat{A}_t^{\mathrm{step}}}{1+w_{\mathrm{step}}}.
\]
最终 scalar advantage 会广播到 action response 的有效 token 上，以复用现有 PPO actor loss。

\section*{实现计划}
\begin{itemize}
\item 新增 \texttt{step\_ppo} estimator，不直接改 GRPO/GiGPO 逻辑。
\item actor backbone 上挂一个 TRL-style value head，不维护独立 critic 模型。
\item rollout 数据保存 \path{uid}、\path{traj_uid}、\path{step_idx}、immediate step reward 和 old scalar state value。
\item return、TD error、GAE 都沿 \path{traj_uid} 与 \path{step_idx} 排序后计算。
\item actor loss 仍复用 token-level PPO clipped surrogate，但 advantage 语义上是 step scalar。
\end{itemize}

\section*{实验矩阵}
\begin{center}
\begin{tabular}{llll}
\toprule
方法 & 是否使用 step 信息 & 是否使用 learned critic & 角色 \\
\midrule
GRPO & 否 & 否 & non-step baseline \\
GiGPO & 是 & 否 & step-aware baseline \\
StepPPO-v1 & 是 & 是 & learned step-aware method \\
StepPPO-v0 & 是 & 是 & normalization ablation \\
\bottomrule
\end{tabular}
\end{center}

\section*{核心指标}
必须记录 validation task score、success rate、text score、valid action ratio、response length、response clipping ratio、KL、policy loss、gradient norm、value loss、value prediction statistics、timing components、GPU memory、CPU memory 和 throughput。

\section*{主要风险}
主要风险包括：value target 高方差、value loss 干扰 actor backbone、token-level PPO ratio 带来的长度偏置、batch shuffle 后 step order 错乱、以及 rollout engine 同步时误包含 value-head 参数。

\end{document}
